% =============================================================================
% 80-lampiran.tex — Lampiran
% =============================================================================
\clearpage
\phantomsection
\addcontentsline{toc}{chapter}{LAMPIRAN}

\chapter*{LAMPIRAN}
\label{chap:lampiran}

% ─────────────────────────────────────────────────────────────────────────────
\section*{Tabel Pemenuhan Kebutuhan/Permintaan Dosen Pemilik Proyek}
\label{lamp:tabel-pemenuhan}

{\small
\begin{longtable}{@{}p{0.6cm} p{4.2cm} p{1.8cm} p{6.4cm}@{}}
  \caption{Pemenuhan Kebutuhan/Permintaan Dosen Pemilik Proyek} \\
  \toprule
  \textbf{No.} & \textbf{Kebutuhan/Permintaan Dosen Pemilik Proyek} & \textbf{Status} & \textbf{Keterangan} \\
  \midrule
  \endfirsthead
  \toprule
  \textbf{No.} & \textbf{Kebutuhan/Permintaan Dosen Pemilik Proyek} & \textbf{Status} & \textbf{Keterangan} \\
  \midrule
  \endhead
  \bottomrule
  \endfoot

  \multicolumn{4}{l}{\textbf{TERCAPAI}} \\
  \midrule

  1 & Penambahan 2 asesmen (GPIUS-2 \& SRS) & \textbf{Tercapai} & Kedua instrumen telah ditambahkan ke katalog dengan butir soal, opsi jawaban, dan rentang interpretasi tervalidasi. \\[0.4em]

  3 & Visualisasi hasil \textit{real-time} & \textbf{Tercapai} & Hasil ditampilkan instan setelah penyelesaian melalui visualisasi grafik (\textit{gauge}, \textit{radar}, diagram batang) menggunakan Recharts. \\[0.4em]

  4 & Permohonan laporan lengkap via surel & \textbf{Tercapai} & Pengguna dapat mengajukan permintaan laporan melalui formulir tervalidasi; permintaan tersimpan pada tabel \texttt{report\_requests} untuk diproses administrator. \\[0.4em]

  5 & Dasbor antrean permintaan (verifikasi/setuju/tolak) & \textbf{Tercapai} & Administrator dapat melihat, menyetujui, atau menolak permintaan laporan; persetujuan memicu generasi PDF dan pengiriman surel. \\[0.4em]

  6 & Visualisasi grafik admin + \textit{filtering} & \textbf{Tercapai} & Dasbor administrator menyajikan visualisasi agregat dengan filter berdasarkan jenis instrumen dan rentang tanggal pada dasbor utama, serta filter lanjutan (nama, jenis kelamin, provinsi, kota, usia, skor, kategori, tanggal) pada halaman data hasil per instrumen. \\[0.4em]

  7 & \textit{Dropdown} domisili bertingkat (Provinsi + Kota/Kab) & \textbf{Tercapai} & Formulir demografis menyediakan pemilihan domisili bertingkat dua level (provinsi $\rightarrow$ kota/kabupaten) menggunakan dataset statis 38 provinsi dari klasifikasi BPS (\textit{Badan Pusat Statistik}). \\[0.4em]

  8 & Opsi `Lainnya' pada Provinsi (manual) & \textbf{Tercapai} & Input provinsi dan kota/kabupaten menyediakan opsi ``Lainnya (ketik manual)'' dengan nilai sentinel \texttt{\_\_other\_\_} yang mengaktifkan kolom teks bebas untuk pengisian wilayah yang tidak terdaftar. \\[0.4em]

  9 & \textit{Save as Image} pada grafik & \textbf{Tercapai} & Komponen grafik dilengkapi tombol unduh gambar menggunakan pustaka \texttt{html-to-image} (fungsi \texttt{toPng}) yang mengkonversi elemen DOM ke format PNG dengan resolusi 2$\times$ \textit{pixel ratio}. \\[0.4em]

  10 & \textit{Dynamic Sorting} (asc/desc) pada tabel & \textbf{Tercapai} & Kolom tabel hasil mendukung pengurutan dinamis \textit{ascending} dan \textit{descending} pada enam parameter (nama, jenis kelamin, usia, skor, kategori, tanggal tes) melalui kueri sisi-server dan pengurutan sisi-klien. \\[0.4em]

  11 & Modul Manajemen Asesmen (CRUD) & \textbf{Tercapai} & Administrator dapat melakukan operasi CRUD penuh atas tes, butir soal, opsi jawaban, dan rentang interpretasi. \\[0.4em]

  12 & Ekspor .CSV dan .XLS & \textbf{Tercapai} & Sistem menyediakan ekspor data dalam format CSV dan XLSX menggunakan pustaka SheetJS. \\[0.4em]

  13 & \textit{Filtered Download} (dinamis) & \textbf{Tercapai} & Data yang diunduh bersifat dinamis mengikuti filter aktif; prosedur \texttt{export} pada \textit{backend} menggunakan fungsi \texttt{buildFilterConditions} yang sama dengan prosedur \texttt{list}, sehingga hasil ekspor mencerminkan tampilan terfilter. \\[0.4em]

  15 & Penyimpanan data sementara (\textit{resume} sesi) & \textbf{Tercapai} & Jawaban tersimpan otomatis pada setiap perubahan melalui prosedur \texttt{saveProgress} dengan pola \textit{upsert} idempoten, memungkinkan pengguna melanjutkan sesi yang terputus. \\[0.4em]

  16 & Modul log aktivitas (\textit{audit trail}) & \textbf{Tercapai} & Setiap modifikasi administratif direkam dalam tabel \texttt{audit\_logs} mencakup aktor, aksi, entitas, serta nilai sebelum dan sesudah. \\[0.4em]

  17 & Pelaporan komprehensif per pengguna & \textbf{Tercapai} & Setiap pengguna memiliki laporan komprehensif mencakup visualisasi data pribadi, analisis per-dimensi, dan interpretasi mendetail. \\[0.4em]

  18 & Ekspor PDF/XLSX/CSV Individual & \textbf{Tercapai} & Hasil mendalam dapat diunduh dalam format PDF (via \texttt{@react-pdf/renderer}), XLSX, dan CSV. \\[0.4em]

  21 & RBAC bertingkat (\textit{super admin}/admin/\textit{researcher}) & \textbf{Tercapai} & Sistem menerapkan RBAC empat peran (\texttt{super\_admin}, \texttt{admin}, \texttt{psychiatrist}, \texttt{researcher}) dengan lima tingkat \textit{middleware} otorisasi. \textit{Super Admin} berakses penuh; Admin tidak dapat mengelola akun; \textit{Psychiatrist} dan \textit{Researcher} memperoleh akses baca saja terhadap data hasil dan statistik. \\[0.4em]

  22 & Modul Manajemen Akun (\textit{Super Admin}) & \textbf{Tercapai} & \textit{Super Admin} dapat membuat, memperbarui, dan mengatur hak akses akun administrator melalui modul \textit{Account Management}. \\[0.4em]

  \midrule
  \multicolumn{4}{l}{\textbf{BELUM TERCAPAI}} \\
  \midrule

  2 & Peringatan otomatis pada pertanyaan belum terisi & \textbf{Belum Tercapai} & Umpan balik eksplisit per-butir belum tercapai karena kompleksitas sinkronisasi status pengisian antara penyimpanan sisi-klien (\texttt{localStorage}) dan persistensi sisi-server secara simultan. Sebagai mitigasi, tombol ``Kirim Asesmen'' dinonaktifkan secara \textit{default} dan hanya aktif setelah seluruh pertanyaan terjawab, sehingga pengiriman tidak lengkap tetap tercegah. \\[0.4em]

  14 & \textit{Import} Data .XLS/.XLSX & \textbf{Belum Tercapai} & Fitur impor data massal belum tercapai karena kompleksitas validasi skema terhadap data historis yang beragam format dan struktur, yang berpotensi menimbulkan inkonsistensi apabila tidak divalidasi secara ketat. Fitur ini ditangguhkan ke iterasi berikutnya demi menjaga integritas data penelitian. \\[0.4em]

  19 & \textit{Contact Us} (pesan ke surel admin) & \textbf{Belum Tercapai} & Fitur korespondensi surel dua arah belum tercapai karena memerlukan integrasi layanan surel dua arah yang melampaui cakupan waktu kerja praktik. Fitur ini diidentifikasi sebagai \textit{backlog} prioritas tinggi untuk fase pengembangan berikutnya. \\[0.4em]

  20 & \textit{Direct Mail Responder} (balasan dari dasbor) & \textbf{Belum Tercapai} & Sama dengan butir 19: memerlukan integrasi surel dua arah yang melampaui cakupan waktu kerja praktik, dan diidentifikasi sebagai \textit{backlog} prioritas tinggi untuk fase berikutnya. \\

\end{longtable}
}

\vspace{2cm}

\noindent Menyetujui, \\
Dosen Pemilik Proyek

\vspace{3cm}

\noindent Dr.\ Yasinta Astin Sokang, S.Psi., M.Psi., Psikolog

\clearpage

% ─────────────────────────────────────────────────────────────────────────────
\section*{Lampiran A: Struktur Direktori Proyek}

Berikut adalah struktur direktori utama proyek Platform Asesmen Digital CHP:

\begin{verbatim}
chpv2/
|-- .github/workflows/ci.yml    # Pipeline CI/CD
|-- src/
|   |-- app/                    # Halaman Next.js (App Router)
|   |   |-- (auth)/             # Grup rute autentikasi
|   |   |-- (protected)/        # Grup rute pengguna terdaftar
|   |   |-- admin/              # Panel administrasi
|   |   |-- test/[slug]/        # Mesin asesmen
|   |   |-- results/[scoreId]/  # Halaman hasil
|   |   +-- tests/              # Katalog tes
|   |-- __tests__/              # 20 file uji (235 kasus uji)
|   |-- lib/
|   |   |-- auth/               # Auth.js v5 (pengguna)
|   |   +-- admin-auth/         # jose JWT (administrator)
|   +-- server/
|       |-- schema/             # 14 file skema Drizzle (20 tabel)
|       |-- scoring/            # 3 file scoring engine
|       +-- trpc/
|           |-- procedures/     # 17 file prosedur (69 prosedur)
|           +-- router.ts       # Komposisi router
|-- drizzle.config.ts
|-- next.config.ts
|-- package.json
|-- tsconfig.json
+-- vitest.config.ts
\end{verbatim}

% ─────────────────────────────────────────────────────────────────────────────
\section*{Lampiran B: Variabel Lingkungan}

Tabel berikut menyajikan seluruh variabel lingkungan yang diperlukan untuk menjalankan platform.

\begin{longtable}{@{}l p{8cm}@{}}
  \toprule
  \textbf{Variabel} & \textbf{Keterangan} \\
  \midrule
  \endfirsthead
  \toprule
  \textbf{Variabel} & \textbf{Keterangan} \\
  \midrule
  \endhead
  \bottomrule
  \endfoot
  \texttt{DATABASE\_URL} & \textit{Connection string} PostgreSQL Neon (\textit{pooled}) \\
  \texttt{AUTH\_SECRET} & Kunci rahasia Auth.js untuk penandatanganan sesi \\
  \texttt{AUTH\_URL} & URL publik aplikasi \\
  \texttt{AUTH\_TRUST\_HOST} & Flag untuk \textit{reverse proxy} \\
  \texttt{ADMIN\_JWT\_SECRET} & Kunci rahasia JWT admin (minimal 32 karakter) \\
  \texttt{ADMIN\_EMAIL} & Surel \textit{bootstrap} \texttt{super\_admin} (opsional) \\
  \texttt{ADMIN\_PASSWORD} & Kata sandi \textit{bootstrap} \texttt{super\_admin} (opsional) \\
  \texttt{UPSTASH\_REDIS\_REST\_URL} & URL REST Upstash Redis \\
  \texttt{UPSTASH\_REDIS\_REST\_TOKEN} & \textit{Token} autentikasi Upstash Redis \\
  \texttt{SENTRY\_DSN} & DSN Sentry untuk pemantauan kesalahan \\
  \texttt{NEXT\_PUBLIC\_SENTRY\_DSN} & DSN Sentry sisi klien \\
  \texttt{SENTRY\_AUTH\_TOKEN} & \textit{Token} autentikasi \textit{build plugin} Sentry \\
  \texttt{RESEND\_API\_KEY} & Kunci API Resend untuk pengiriman surel \\
  \texttt{ENCRYPTION\_KEY} & Kunci AES-256-GCM (64 karakter heksadesimal) \\
  \texttt{GOOGLE\_CLIENT\_ID} & ID klien Google OAuth (belum diaktifkan) \\
  \texttt{GOOGLE\_CLIENT\_SECRET} & Rahasia klien Google OAuth (belum diaktifkan) \\
\end{longtable}

% ─────────────────────────────────────────────────────────────────────────────
\section*{Lampiran C: Instruksi \textit{Setup} Lokal}

Untuk menjalankan proyek secara lokal, ikuti langkah-langkah berikut:

\begin{enumerate}[leftmargin=2cm]
  \item \textit{Clone} repositori: \texttt{git clone <repository-url>}
  \item Instal dependensi: \texttt{pnpm install}
  \item Salin \texttt{.env.example} ke \texttt{.env.local} dan isi seluruh variabel
  \item Jalankan migrasi basis data: \texttt{pnpm drizzle-kit push}
  \item Jalankan \textit{seed} data: \texttt{node run-drizzle.js}
  \item Jalankan \textit{development server}: \texttt{pnpm dev}
  \item Akses platform di \texttt{http://localhost:3000}
\end{enumerate}

Target waktu \textit{setup}: $\leq$ 15 menit untuk pengembang baru (Tujuan KP \#6).
